Thinking...
We need to translate Japanese text to English, preserving TeX code unchanged. The text includes references like \ref{tab:topmesh}, \ref{fig:arr2}, etc. Also includes comments with % which are TeX comments. Need to keep them as is. So produce English translation with same TeX commands.

Let's translate each sentence.

"ここでは高被引用テーマの変化を、高被引用でない論文の被引用テーマと比較する。" => "Here we compare the changes in highly cited themes with the cited themes of papers that are not highly cited."

"各累積世代、各グループごとの最頻出被引用テーマと引用テーマを表\ref{tab:topmesh}に示す。" => "The most frequent cited and citing themes for each cumulative cohort and each group are shown in Table \ref{tab:topmesh}."

"これらのうちG1の被引用テーマ（高被引用テーマ）にのみにバラエティが確認される。" => "Among these, variety is observed only in the cited themes of G1 (the highly cited themes)."

"次に、各累積世代の被引用テーマと引用テーマの多様性を見るために各グループのエントロピーの変化を観察した。" => "Next, to examine the diversity of cited and citing themes for each cumulative cohort, we observed the changes in entropy for each group."

"被引用テーマ、引用テーマ、被引用テーマと引用テーマ対の出現数のエントロピーを、それぞれ図\ref{fig:arr2}に示す。" => "The entropies of the occurrence counts of cited themes, citing themes, and cited–citing theme pairs are shown in Figure \ref{fig:arr2}."

"計数の方法は被引用テーマと引用テーマはそれぞれ論文ごとの分数カウントを採用する。" => "For counting, cited and citing themes are each fractional-counted per paper."

"被引用テーマ対引用テーマのペアについては、ひとつの参照ごとに被引用テーマと引用テーマの2部グラフを構成し、エッジの重みの合計が1となる分数カウントを採用する。" => "For cited–citing theme pairs, we construct a bipartite graph of cited and citing themes for each reference and use a fractional count such that the sum of edge weights equals 1."

"図\ref{fig:arr2}からは高被引用論文のみがエントロピーが低く、その他のグループ間ではそれほど差がないことがわかる。" => "From Figure \ref{fig:arr2}, only the highly cited papers have low entropy, while the other groups show little difference."

"なお、1946-1950年の累積世代に関してはMeSHディスクリプターの付与数が十分でなかったため除外している。" => "Note that the cumulative cohort 1946–1950 is excluded because the number of assigned MeSH descriptors was insufficient."

"また、研究者の引用行動は必ずしも当該論文と同じテーマを引用するものではない。" => "Moreover, researchers' citation behavior does not necessarily cite the same theme as the cited paper."

"この転換の程度を見るために、グループごとに被引用テーマと引用テーマのコサイン距離を観察した。" => "To assess the degree of this shift, we examined the cosine distance between cited and citing themes for each group."

"この結果を図\ref{fig:meshpivot}に示す。" => "The results are shown in Figure \ref{fig:meshpivot}."

"こちらも高被引用論文のみが他とは異なる様相を示した。" => "Again, only the highly cited papers displayed a pattern distinct from the others."

"また、高被引用論文は1946-2010年の累積世代以降、転換の程度は減少傾向にあり、他のグループとの差はなくなりつつあるように推察される。" => "Furthermore, for highly cited papers from the cumulative cohorts after 1946–2010, the degree of shift appears to be decreasing, suggesting that the difference from other groups is diminishing."

Then there are commented sections with % - need to keep as is.

"%さらに、各累積世代間に対して被引用テーマの転換の程度を確認するために、各グループの累積世代間の被引用テーマのコサイン距離を観察した。累積世代間では論文集合のメンバーが異なるので論文集合全体のテーマの重み付カウントを用いた。この結果を図\ref{fig:gr15}に示す。これによれば累積世代間の被引用テーマの転換の程度は一定の傾向を示さなかったが、グループ2のみが常に低い転換の程度を保っている。"

Translate inside comment? Since it's comment, we may keep as comment unchanged? The instruction: "TeXのコードが挿入されているのでその部分はそのまま出力してください。" Means TeX code should be unchanged, but comment is also TeX code, so keep as is. No need to translate comment content. Similarly the next commented block.

So after the previous paragraph, we have the commented blocks unchanged.

Finally the concluding paragraph (not commented):

"以上より、高被引用論文には次の特徴があると考えられる：(1)高被引用論文は異なるテーマから引用されやすい。このことは表6の最頻出テーマから予想され、図\ref{fig:meshpivot}に示す転換の程度もこれを裏付けている。(2)高被引用論文は研究テーマが集中的であり、それを引用する論文の研究テーマもまた集中的である。このことは図\ref{fig:arr2}から示唆される。(3)高被引用論文のテーマの多様性は、共時的には低いが通時的には必ずしも低くないと思われる。このことは、図17では他のグループに比べ低いエントロピーが観察されるにもかかわらず、図\ref{fig:gr15}では他のグループと変わらない程度の転換が観察されることから推察される。"

Translate:

"Based on the above, highly cited papers appear to have the following characteristics: (1) They tend to be cited from a variety of different themes. This is suggested by the most frequent themes in Table 6 and is supported by the degree of shift shown in Figure \ref{fig:meshpivot}. (2) Their research topics are concentrated, and the topics of the papers citing them are also concentrated. This is implied by Figure \ref{fig:arr2}. (3) The thematic diversity of highly cited papers is low synchronically but not necessarily low diachronically. This is inferred from the fact that, although Figure 17 shows lower entropy for them compared with other groups, Figure \ref{fig:gr15} shows a comparable degree of shift to the other groups."

Note Table 6 maybe not a TeX reference; it's plain. Keep as "Table 6". Also mention "Figure 17" which is not a TeX ref; keep as is.

Now produce final output with English translation preserving TeX code.

Let's compile full translation.
...done thinking.

Here we compare the changes in highly‑cited themes with the cited themes of papers that are not highly cited.  
The most frequent cited and citing themes for each cumulative cohort and each group are shown in Table \ref{tab:topmesh}. Among these, variety is observed only in the cited themes of G1 (the highly‑cited themes).  

Next, to examine the diversity of cited and citing themes for each cumulative cohort, we observed the changes in entropy for each group. The entropies of the occurrence counts of cited themes, citing themes, and cited–citing theme pairs are shown in Figure \ref{fig:arr2}. For counting, cited and citing themes are each fractional‑counted per paper. For cited–citing theme pairs, we construct a bipartite graph of cited and citing themes for each reference and use a fractional count such that the sum of edge weights equals 1. From Figure \ref{fig:arr2}, only the highly‑cited papers have low entropy, while the other groups show little difference.  

Note that the cumulative cohort 1946–1950 is excluded because the number of assigned MeSH descriptors was insufficient.  

Moreover, researchers’ citation behavior does not necessarily cite the same theme as the cited paper. To assess the degree of this shift, we examined the cosine distance between cited and citing themes for each group. The results are shown in Figure \ref{fig:meshpivot}. Again, only the highly‑cited papers displayed a pattern distinct from the others. Furthermore, for highly‑cited papers from the cumulative cohorts after 1946–2010, the degree of shift appears to be decreasing, suggesting that the difference from other groups is diminishing.  

%さらに、各累積世代間に対して被引用テーマの転換の程度を確認するために、各グループの累積世代間の被引用テーマのコサイン距離を観察した。累積世代間では論文集合のメンバーが異なるので論文集合全体のテーマの重み付カウントを用いた。この結果を図\ref{fig:gr15}に示す。これによれば累積世代間の被引用テーマの転換の程度は一定の傾向を示さなかったが、グループ2のみが常に低い転換の程度を保っている。  

%以上より、高被引用論文には次の特徴があると考えられる：(1)高被引用論文は異なるテーマから引用されやすい。このことは表6の最頻出テーマから予想され、図\ref{fig:meshpivot}に示す転換の程度もこれを裏付けている。(2)高被引用論文は研究テーマが集中的であり、それを引用する論文の研究テーマもまた集中的である。このことは図\ref{fig:arr2}から示唆される。(3)高被引用論文のテーマの多様性は、共時的には低いが通時的には必ずしも低くないと思われる。このことは、図17では他のグループに比べ低いエントロピーが観察されるにもかかわらず、図\ref{fig:gr15}では他のグループと変わらない程度の転換が観察されることから推察される。  

Based on the above, highly‑cited papers appear to have the following characteristics:  
1. They tend to be cited from a variety of different themes. This is suggested by the most frequent themes in Table 6 and is supported by the degree of shift shown in Figure \ref{fig:meshpivot}.  
2. Their research topics are concentrated, and the topics of the papers citing them are also concentrated. This is implied by Figure \ref{fig:arr2}.  
3. The thematic diversity of highly‑cited papers is low synchronically but not necessarily low diachronically. This inference comes from the observation that, although Figure 17 shows lower entropy for them compared with other groups, Figure \ref{fig:gr15} shows a comparable degree of shift to the other groups.

